DGFiP, REI 2024 et DVF 2023-2024 ; Insee, recensement 2022, comptes de patrimoine 2024 et Insee Analyses n° 91 ; calculs de l'auteur (synthese/calculs\_foncier.py, foncier/simulation.py).\par Note : réforme naïve : droits de mutation sur les logements divisés par deux (5,7 Md\euro{}), compensés par une hausse de 24 \% de la taxe foncière actuelle dans toutes les communes, sans réaction des ventes. La baisse des droits est rapportée à l'ensemble des logements de la commune : c'est ce que paie en moyenne chaque année un logement au fil des ventes. Un ménage qui ne vend ni n'achète paie la seule hausse de taxe foncière (barres grises).
